\subsection{Semplificazioni algebriche}
\label{sec:esercizi-limiti-semplificazioni-algebriche}

\begin{exercise}[Fattorizzazione di una forma indeterminata]
\label{ex:limiti-foglio-2311-04}
\difficoltaesercizio{0}
Calcolare il limite
\[
  \lim_{x\to2}\frac{x^2-3x+2}{x^2+2-6}.
\]
\rimandoteoria{thm:lez14-limite-rapporto}
\end{exercise}
\begin{proof}[Soluzione dell'esercizio]
Per ogni $x$ nel dominio, fattorizzando e semplificando si ottiene
\[
  \frac{x^2-3x+2}{x^2+2-6}
  =\frac{(x-1)(x-2)}{(x-2)(x+2)}
  =\frac{x-1}{x+2}.
\]
Pertanto
\[
  \lim_{x\to2}\frac{x^2-3x+2}{x^2+2-6}
  =\lim_{x\to2}\frac{x-1}{x+2}
  =\frac14.
\]
\end{proof}

\begin{exercise}[Semplificazione di un fattore comune]
\label{ex:limiti-foglio-2311-05}
\difficoltaesercizio{0}
Calcolare il limite
\[
  \lim_{x\to0}\frac{x^2+3x}{x}.
\]
\rimandoteoria{thm:lez14-limite-somma}
\end{exercise}
\begin{proof}[Soluzione dell'esercizio]
Per ogni $x$ nel dominio, raccogliendo $x$ e semplificando si ottiene
\[
  \frac{x^2+3x}{x}=x+3.
\]
Pertanto
\[
  \lim_{x\to0}\frac{x^2+3x}{x}
  =\lim_{x\to0}(x+3)=3.
\]
\end{proof}

\begin{exercise}[Semplificazione di due trinomi]
\label{ex:limiti-foglio-2311-06}
\difficoltaesercizio{0}
Calcolare il limite
\[
  \lim_{x\to3}\frac{x^2-8x+15}{x^2-10x+21}.
\]
\rimandoteoria{thm:lez14-limite-rapporto}
\end{exercise}
\begin{proof}[Soluzione dell'esercizio]
Per ogni $x$ nel dominio, fattorizzando e semplificando si ottiene
\[
  \frac{x^2-8x+15}{x^2-10x+21}
  =\frac{(x-3)(x-5)}{(x-3)(x-7)}
  =\frac{x-5}{x-7}.
\]
Pertanto
\[
  \lim_{x\to3}\frac{x^2-8x+15}{x^2-10x+21}
  =\lim_{x\to3}\frac{x-5}{x-7}
  =\frac{-2}{-4}=\frac12.
\]
\end{proof}

\clearpage

\begin{exercise}[Semplificazione di una differenza di quadrati]
\label{ex:limiti-foglio-2311-08}
\difficoltaesercizio{0}
Calcolare il limite
\[
  \lim_{x\to2}\frac{x^2-4}{x^2-3x+2}.
\]
\rimandoteoria{thm:lez14-limite-rapporto}
\end{exercise}
\begin{proof}[Soluzione dell'esercizio]
Per ogni $x$ nel dominio, fattorizzando e semplificando si ottiene
\[
  \frac{x^2-4}{x^2-3x+2}
  =\frac{(x-2)(x+2)}{(x-2)(x-1)}
  =\frac{x+2}{x-1}.
\]
Pertanto
\[
  \lim_{x\to2}\frac{x^2-4}{x^2-3x+2}
  =\lim_{x\to2}\frac{x+2}{x-1}=4.
\]
\end{proof}
